\section{Decision log}\label{app:decisions}

Each entry records a choice that could affect the results or the usability of the project, the alternatives that were considered and the reason for the choice.
Entries are never deleted: a decision that is reversed is superseded by a new entry that refers to it.

\decision{D1}{18 October 2026}{Language of the project}
{The paper, the code and its comments are written in English.}
{French.}
{English is the working language of trading floors and of open-source code, and the project is meant to be shown to recruiters, including outside France.}

\decision{D2}{18 October 2026}{Code organised as an installable Python package}
{The code lives in a package (\texttt{src/fxlab}) installed in a virtual environment; notebooks only import and use it.}
{(a) Code written directly in Jupyter notebooks. (b) Loose scripts.}
{Code in notebooks cannot be tested or reused, and is easily run out of order. A package separates the library (tested, reused by every phase) from the analysis (notebooks). The \texttt{src} layout guarantees that the tests run against the installed package, not against stray files.}

\decision{D3}{18 October 2026}{Course examples used as automated tests}
{Every worked example of the course that the code can reproduce is encoded as a test (\texttt{pytest}), with a tolerance that reflects the rounding of the published inputs.}
{Checking results by eye in notebooks.}
{The course examples have known answers, so they are an independent specification of what the code must do. Automated tests are rerun after every change, which prevents an improvement in one place from silently breaking another.}

\decision{D4}{18 October 2026}{Valuation by propagation of amounts}
{Portfolios are valued by propagating the amount held in each constituent and resetting it on rebalancing dates (course, section~1.8).}
{Compounding the portfolio return $\w'\vect{r}$ with fixed weights.}
{Compounding with fixed weights is only valid when the portfolio is rebalanced at every date. Propagating amounts handles any rebalancing frequency and gives the amount traded, which is needed for transaction costs.}

\decision{D5}{18 October 2026}{Free, official data sources}
{Exchange rates and interest rates are taken from official publishers (Federal Reserve, OECD, ECB) through FRED and the ECB data portal; forward rates are implied by CIP.}
{(a) Bloomberg or Refinitiv: no access. (b) Yahoo Finance: unofficial interface, no documentation of the series, known data errors.}
{Official sources are documented, stable and allow anyone to reproduce the results. The cost is that forward rates are not observed; this is acknowledged as a limitation (Section~\ref{sec:data}).}

\decision{D6}{18 October 2026}{Working paper in \LaTeX{}, versioned with the code}
{The documentation is a \LaTeX{} article stored in the same repository as the code.}
{(a) A word processor document. (b) Notebooks only.}
{\LaTeX{} is the standard for scientific documents with equations. Keeping it in the repository ties each version of the text to the version of the code that produced its results.}

\decision{D7}{19 October 2026}{Data downloaded by a script from the FRED CSV interface}
{A script downloads each series from its public FRED address and caches the raw file in \texttt{data/raw/}; the files are read and cleaned by the \texttt{fxlab.data} module. Spot rates are converted to US dollars per unit of foreign currency.}
{(a) Manual download from the website. (b) The FRED API, which requires a personal key. (c) A third-party package such as \texttt{pandas-datareader}.}
{A script makes the dataset reproducible with one command and can refresh it at any time. The public CSV interface needs no key and only the Python standard library, so there is no dependency that could break. A single quotation convention avoids sign errors on half of the currencies.}

\decision{D8}{19 October 2026}{Monthly frequency, month-end sampling}
{Currency strategies are formed and evaluated monthly, using the last available spot rate of each month.}
{(a) Daily frequency. (b) Weekly frequency. (c) Monthly averages of spot rates.}
{Interest rates are only available monthly, so a higher frequency would reuse the same rate for a whole month. Monthly rebalancing is the standard of the currency literature \citep{lustig2011,menkhoff2012carry}, which makes the results comparable, and it keeps turnover and transaction costs low. Month-end rates, not averages, are used because a position is opened and closed at a given rate, not at an average.}

\decision{D9}{19 October 2026}{Unbalanced panel from January 1999}
{The sample starts in January 1999; each currency enters as soon as its spot rate and interest rate are both available. A balanced panel (2002--2026) is used as a robustness check.}
{(a) Balanced panel from May 2002 only. (b) Start in the 1980s without the euro.}
{Starting with the euro gives a homogeneous G10 universe. An unbalanced panel uses three more years of data without inventing any, as in \citet{lustig2011}; the robustness check verifies that the conclusions do not depend on this choice.}

\decision{D10}{19 October 2026}{US dollar as base currency for the currency portfolios}
{The excess returns of Phase 2 are measured from the point of view of a US investor; the euro-based investor is the subject of Phase 3.}
{Euro as base currency from the start.}
{The data are quoted against the US dollar and almost all of the literature uses a US investor, so the results can be compared with published ones. Excess returns of long-short portfolios do not depend much on the base currency, since the base currency cancels out of a long-short position.}

\decision{D11}{19 October 2026}{Excess return defined as the simple return of a one-month forward}
{$rx_{t+1} = S_{t+1}/F_t - 1$ with $F_t$ implied by CIP and the monthly rate equal to the annual rate divided by 12 (equations~\ref{eq:fwd-monthly}--\ref{eq:rx-monthly}).}
{(a) Log excess return $\ln S_{t+1} - \ln F_t$, common in the literature. (b) Deposit version \eqref{eq:rx-deposit}.}
{The portfolio formula $r_p = \w'\vect{r}$ of the course holds for simple returns, not for log returns, so simple returns can be aggregated across currencies without approximation. The forward is the instrument actually traded by an FX desk for this exposure, and it requires no capital beyond margin. The deposit version differs only by the factor $1+i_t$ (equation~\ref{eq:rx-forward}).}

\decision{D12}{19 October 2026}{Continuous integration and MIT licence}
{The test suite runs automatically on GitHub Actions at every push, on Python 3.11 and 3.13; the code is published under the MIT licence.}
{(a) Running the tests locally only. (b) No licence, or a restrictive one.}
{Tests run on a clean machine catch errors that a local environment hides (a missing dependency, a file that only exists locally), and the badge shows any reader that the suite passes. Without a licence, published code cannot legally be reused; MIT is the most common licence for demonstration projects and imposes only attribution.}

\decision{D13}{19 October 2026}{Compiled paper versioned in the repository}
{The compiled PDF is copied to \texttt{paper/fx-portfolio-lab.pdf} by \texttt{make} and committed with each update of the paper; LaTeX build files remain ignored.}
{(a) PDF attached to GitHub releases only. (b) PDF compiled by the continuous integration. (c) Sources only.}
{A reader who is not going to compile \LaTeX{} must be able to open the paper in one click from the README. Committing the PDF is the simplest way to guarantee it and keeps each version of the PDF next to the code that produced it. Compiling in the CI would require installing a TeX distribution on every run, which is slow for a document that changes only a few times per phase.}

\decision{D14}{19 October 2026}{Transaction costs: roll plus change of position}
{The notional kept from one month to the next pays the swap half-spread (roll); the notional traded pays the outright half-spread (equation~\ref{eq:cost}).}
{A single half-spread on the notional traded only.}
{Holding a forward position requires rolling it every month with an FX swap, which is how an FX desk maintains the exposure. Ignoring the roll would understate the cost of strategies that hold their positions; on the dollar factor, the roll accounts for most of the cost (Section~\ref{sec:dollar}).}

\decision{D15}{19 October 2026}{Cost level: stated assumption, sensitivity and break-even}
{The base-case half-spreads are an explicit assumption by liquidity tier (Section~\ref{sec:costs}); results are reported for the base case scaled by 0, 1, 2 and 5, with the break-even cost scale of each strategy.}
{(a) Figures presented as estimates without a public data source. (b) A single cost for all currencies.}
{Bid--ask quotes are not freely available, so any single figure would be an assumption; presenting it as an estimate would overstate its precision. The break-even scale answers the question that matters (up to which cost does a strategy remain profitable?) without depending on the assumption, and the tiers keep the ranking of liquidity between currencies.}

\decision{D16}{20 October 2026}{Interbank rates in covered interest parity, with a risk-free-rate robustness check}
{Forwards are implied from 3-month interbank rates over the whole sample; from November 2019, euro and sterling excess returns are recomputed with the overnight risk-free rates (SOFR, ESTR, SONIA) as a robustness check (Section~\ref{sec:rates-choice}).}
{(a) Risk-free rates only: they do not exist before 2018--2019 and not for every currency in free data. (b) Interbank rates without discussion.}
{Interbank rates contain a bank credit premium and a term component that the rates of a collateralised FX swap do not contain, a point raised by practitioners. But they were the market reference before the benchmark reform, they cover the whole sample and all currencies, and the robustness check shows that switching to risk-free rates changes monthly excess returns by about 1 basis point. The choice is therefore kept and its limits are stated.}

\decision{D17}{20 October 2026}{Extension of the euro and sterling rates after January 2026}
{The euro rate is continued with the 3-month Euribor published by the ECB; the sterling rate with the monthly average of SONIA, shifted by the mean gap over the last twelve months of overlap.}
{(a) Let the two currencies leave the panel in February 2026. (b) Splice SONIA without a shift.}
{The OECD euro series is the Euribor, so the ECB series continues it exactly. For sterling no 3-month interbank rate exists after LIBOR; a constant shift removes the jump at the junction ($-0.08$ point) without changing the monthly variation of the rate. Keeping the euro and sterling in the panel avoids changing the investment universe at the end of the sample.}

\decision{D18}{20 October 2026}{Euro drivers measured with monthly averages and Newey--West standard errors}
{The regressions of Section~\ref{sec:euro} use changes of monthly averages for the exchange rates as well as for rates and yields, and Newey--West standard errors with the usual lag rule.}
{(a) Month-end exchange rates against monthly-average yields. (b) Classical OLS standard errors.}
{The OECD rates and yields are only available as monthly averages; mixing them with month-end exchange rates would misalign the timing of the two sides. Averaging induces autocorrelation in monthly changes, and financial data are heteroskedastic, so classical standard errors would overstate precision; Newey--West standard errors are robust to both.}

\decision{D19}{20 October 2026}{Scope: what is not added from the practitioner advice}
{ECB meeting pricing, sovereign debt auctions and the carry and roll-down of bonds are not part of the project; they are followed as market monitoring outside it.}
{(a) An event study of ECB meeting days. (b) A rates module for bond carry, roll-down and repo.}
{Measuring a monetary policy surprise requires the market pricing of each meeting (dated overnight-index swaps), which is not freely available historically, and a complete list of meeting dates that would have to be compiled by hand; auctions have no free historical database. Bond carry and roll-down belong to a rates desk: adding them would blur a project whose subject is currencies. The euro is studied instead through the two channels of Section~\ref{sec:euro}, which public data can measure.}

\subsection*{Open decisions}
\addcontentsline{toc}{subsection}{Open decisions}
None at this stage. The decisions of Phase 2 (signal construction, portfolio formation, estimation of the covariance matrix) will be added before the corresponding code is written.
